\chapter*{Abstract}
\addcontentsline{toc}{chapter}{ABSTRACT} % to add abstract to table of contents
\begin{spacing}{1.5}
Integrity is the practice of being honest and showing a consistent and uncompromising adherence to strong moral and ethical principles and values. In ethics, integrity is regarded as the honesty and truthfulness or accuracy of one's actions. Integrity can stand in opposition to hypocrisy, in that judging with the standards of integrity involves regarding internal consistency as a virtue, and suggests that parties holding within themselves apparently conflicting values should account for the discrepancy or alter their beliefs. The word integrity evolved from the Latin adjective integer, meaning whole or complete. In this context, integrity is the inner sense of ``wholeness" deriving from qualities such as honesty and consistency of character. As such, one may judge that others ``have integrity" to the extent that they act according to the values, beliefs and principles they claim to hold.
\paragraph{}
%\noindent uncomment this command if tab not required
An individual's value system provides a framework within which the individual acts in ways which are consistent and expected. Integrity can be seen as the state or condition of having such a framework, and acting congruently within the given framework.One essential aspect of a consistent framework is its avoidance of any unwarranted (arbitrary) exceptions for a particular person or group—especially the person or group that holds the framework. In law, this principle of universal application requires that even those in positions of official power be subject to the same laws as pertain to their fellow citizens.
\end{spacing}

